\section{Row--orbit energy and common-target pruning}
\label{sec:tpc30-row-orbit}

For $c\ge1$, define the absolute row congruence mass
\begin{equation}
 \cR_c
 :=\sum_{\alpha,\gamma}
 |\gamma_\alpha^{(1)}\gamma_\gamma^{(2)}|
 \one_{c\mid m_\alpha-m_\gamma}.
 \label{eq:tpc30-row-congruence-mass}
\end{equation}
The following estimate was used for reflected content in TPC-29.  We
record it in the present notation because it is now paired with the
full target gcd rather than a selected divisor gcd.

\begin{lemma}[Opened-row residue energy]
\label{lem:tpc30-row-residue-energy}
Uniformly for $c\ge1$,
\begin{equation}
 \boxed{
 \cR_c
 \ll_{\mathscr D}
 \left(\frac{Q^2}{c}+Q\right)\log(2L).}
 \label{eq:tpc30-row-residue-bound}
\end{equation}
\end{lemma}

\begin{proof}
For $i=1,2$ and $a\pmod c$, put
\[
 X_i(a)=\sum_{m_\alpha\equiv a\, (c)}
 |\gamma_\alpha^{(i)}|,
 \qquad
 E_i(a)=\sum_{m_\alpha\equiv a\, (c)}
 |\gamma_\alpha^{(i)}|^2.
\]
The row integers are distinct and lie in an interval of length
$O_{\mathscr D}(Q)$.  Each residue class contains
$O_{\mathscr D}(Q/c+1)$ rows, so Cauchy--Schwarz gives
\[
 X_i(a)^2
 \ll_{\mathscr D}\left(\frac Qc+1\right)E_i(a).
\]
It follows that
\begin{align*}
 \cR_c
 &=\sum_{a\, (c)}X_1(a)X_2(a)\\
 &\ll_{\mathscr D}
 \left(\frac Qc+1\right)
 \left(\sum_aE_1(a)\right)^{1/2}
 \left(\sum_aE_2(a)\right)^{1/2}.
\end{align*}
Apply the two estimates in \eqref{eq:tpc30-row-l2}.
\end{proof}

For a row pair and orbit point, abbreviate
\begin{equation}
 \cG_{\alpha,\gamma}(j)
 :=(N_\alpha(j),N_\gamma(j)).
 \label{eq:tpc30-physical-target-gcd}
\end{equation}
Taking absolute values in any raw channel gives
\begin{equation}
 |\cK^{(\nu)}_{\alpha,\gamma}(j)|
 \ll_{\eps,\mathscr D}X^\eps,
 \qquad 1\le\nu\le3,
 \label{eq:tpc30-kernel-envelope}
\end{equation}
by \eqref{eq:tpc30-pointwise-envelope}.

\begin{proposition}[A uniform exact-content layer]
\label{prop:tpc30-exact-content-layer}
For every integer $c\ge1$ and $1\le\nu\le3$,
\begin{equation}
 \boxed{
 |\cS_\nu(\cG_{\alpha,\gamma}(j)=c)|
 \ll_{\eps,h,\mathscr D}X^\eps
 \left(1+\frac Jc\right)
 \left(\frac{Q^2}{c}+Q\right).}
 \label{eq:tpc30-exact-content-layer}
\end{equation}
\end{proposition}

\begin{proof}
The mask in \eqref{eq:tpc30-row-mask} first deletes the diagonal rows.
By \cref{lem:tpc30-target-gcd}, the exact layer is empty unless
$c\mid m_\alpha-m_\gamma$.  For each surviving row pair,
\cref{lem:tpc30-content-occupancy} gives
$O_{\mathscr D}(1+J/c)$ orbit points.  Apply
\eqref{eq:tpc30-kernel-envelope}, sum the row weights through
$\cR_c$, and use \eqref{eq:tpc30-row-residue-bound}.
\end{proof}

Summing \eqref{eq:tpc30-exact-content-layer} over every integer
$c>Z$ would introduce an avoidable endpoint term.  The stronger
tail estimate comes from first fixing the row pair.

\begin{lemma}[Fixed-row large-content occupancy]
\label{lem:tpc30-fixed-row-large-content}
Let $m\ne n$ be physical rows, let $\cI$ be an interval of
$O_{\mathscr D}(J)$ consecutive integers, and let $Z\ge1$.  Then
\begin{equation}
 \boxed{
 \#\{j\in\cI:(j,h)=1,\ \cG_{m,n}(j)>Z\}
 \ll_{\eps,h,\mathscr D}
 X^\eps\left(1+\frac JZ\right).}
 \label{eq:tpc30-fixed-row-large-content}
\end{equation}
\end{lemma}

\begin{proof}
Put $\Delta=|m-n|\ne0$.  Every exact value
$c=\cG_{m,n}(j)$ divides $\Delta$ by
\cref{lem:tpc30-target-gcd}.  Grouping by the exact value and applying
\cref{lem:tpc30-content-occupancy} gives
\begin{align}
 \#\{j\in\cI:(j,h)=1,\ \cG_{m,n}(j)>Z\}
 &\le
 \sum_{\substack{c\mid\Delta\\c>Z}}
 O_{\mathscr D}\left(1+\frac Jc\right)
 \notag\\
 &\ll_{\mathscr D}
 \tau(\Delta)\left(1+\frac JZ\right).
 \label{eq:tpc30-fixed-row-divisor-sum}
\end{align}
Since $\Delta\ll_{\mathscr D}Q\ll X^{O_{\mathscr D}(1)}$, the standard
divisor bound absorbs $\tau(\Delta)$ into $X^\eps$.
\end{proof}

\begin{theorem}[Common-target-content pruning]
\label{thm:tpc30-common-target-pruning}
Let $Z\ge1$.  For each $1\le\nu\le3$, the physical raw channel in
\eqref{eq:tpc30-restricted-raw-channel}, restricted by
\begin{equation}
 \cE_Z(\alpha,\gamma,j)
 :=\{\cG_{\alpha,\gamma}(j)>Z\},
 \label{eq:tpc30-large-target-content-event}
\end{equation}
satisfies
\begin{equation}
 \boxed{
 |\cS_\nu(\cE_Z)|
 \ll_{\eps,h,\mathscr D}X^\eps Q^2
 \left(1+\frac JZ\right)
 \ll_{\eps,h,\mathscr D}X^\eps
 \left(Q^2+\frac{XQ}{Z}\right).}
 \label{eq:tpc30-main-content-bound}
\end{equation}
The same bound holds for the sum of the three raw channels.
\end{theorem}

\begin{proof}
The off-diagonal mask gives $m_\alpha\ne m_\gamma$, so
\cref{lem:tpc30-fixed-row-large-content} applies to every row pair.
Use \eqref{eq:tpc30-kernel-envelope} and then the row $\ell^1$ bounds:
\begin{align*}
 |\cS_\nu(\cE_Z)|
 &\ll X^\eps\left(1+\frac JZ\right)
 \sum_{\alpha,\gamma}
 |\gamma_\alpha^{(1)}\gamma_\gamma^{(2)}|\\
 &\ll X^\eps Q^2\left(1+\frac JZ\right).
\end{align*}
Finally use $JQ\asymp X$.  If $Z\gg_{\mathscr D}Q$, the event is
empty because
\begin{equation}
 \cG_{\alpha,\gamma}(j)
 \mid m_\alpha-m_\gamma
 \ll_{\mathscr D}Q.
 \label{eq:tpc30-content-upper-range}
\end{equation}
\end{proof}

\begin{corollary}[Power-saving form]
\label{cor:tpc30-content-power-saving}
Suppose \eqref{eq:tpc30-row-scales} holds and $Z\ge X^\eta$ for a
fixed $\eta>0$.  Then, for every fixed
\begin{equation}
 \eta_0<\min\{\eta,1-\beta\},
 \label{eq:tpc30-saving-exponent}
\end{equation}
one has
\begin{equation}
 \boxed{
 \sum_{\nu=1}^3|\cS_\nu(\cE_Z)|
 \ll_{\eta,\eta_0,h,\mathscr D}XQX^{-\eta_0}.}
 \label{eq:tpc30-content-power-saving}
\end{equation}
\end{corollary}

\begin{proof}
Relative to $XQ\asymp JQ^2$, the two terms in
\eqref{eq:tpc30-main-content-bound} have respective ratios
$J^{-1}$ and $Z^{-1}$.  Use \eqref{eq:tpc30-row-scales}; the strict
gap in \eqref{eq:tpc30-saving-exponent} absorbs logarithms and the
$o(1)$ scale factors.
\end{proof}

\begin{remark}[The absolute one-point floor]
\label{rem:tpc30-one-point-floor}
The term $Q^2$ comes from the $+1$ in the occupancy of each exact
content class.  It is power-small because $J$ is a fixed power, but
the present argument does not replace it by $JQ$.  This is a
boundary of the absolute fixed-row proof, not an arithmetic lower
bound or an optimality statement.
\end{remark}
